% 原书物理页=印刷页16--19；1.6 全节.
% 注：1.6 节的结尾（小结之后的段落、1.6.7 命题及其证明、习题）印于原书物理页 20，属 1.6 节内容.
\section{$\mathrm{GL}_n(\mathbb R)$ 的闭子群}\label{sec:lie-1-6}
我们此处的目标是证明 $\mathrm{GL}_n(\mathbb R)$ 的任何闭子群都典范地是一个 Lie 子群（Cartan--Von Neumann 定理的特例），并给出一种计算其 Lie 代数的方法. 因此，这将应用于上文介绍的经典群. 我们的主要工具将是\newterm{lie-matrix-exponential}{矩阵指数映射}{matrix exponential}.
\begin{numpara}[指数映射]\label{np:lie-1-6-1}
设 $V$ 是一个有限维 $\mathbb R$-向量空间. 在 $V$ 上取一个范数，并在 $\operatorname{End}_{\mathbb R}(V)$ 上配备在单位球上取上确界所得的范数 $\|.\|$. 这一范数是次乘性的：$\|uv\|\leqslant\|u\|.\|v\|$. 由此可知，级数
\[
\exp(u):=\sum_{n\in\mathbb N}\frac{1}{n!}u^n
\]
是正规收敛的，从而定义了一个解析映射，因而是从 $\operatorname{End}_{\mathbb R}(V)$ 到其自身的光滑映射. 若 $u$ 与 $v$ 交换，则有 $\exp(u+v)=\exp(u)\exp(v)$. 特别地，$\exp(u)$ 可逆，其逆为 $\exp(-u)$. 我们有 $\exp(0)=\operatorname{id}_V$，而 $d_0\exp=\operatorname{id}_{\operatorname{End}_{\mathbb R}(V)}$ 可逆. 于是，$\exp$ 实现了从 $\operatorname{End}_{\mathbb R}(V)$ 中 $0$ 的一个邻域到 $\mathrm{GL}(V)$ 中 $\operatorname{id}_V$ 的一个邻域的微分同胚.
在以上整个讨论中，当然可以把 $\operatorname{End}_{\mathbb R}(V)$ 换成 $\mathrm{M}_n(\mathbb R)$，把 $\mathrm{GL}(V)$ 换成 $\mathrm{GL}_n(\mathbb R)$.
\end{numpara}
\begin{numpara}[$\mathrm{GL}_n(\mathbb R)$ 的单参数子群]\label{np:lie-1-6-2}
指数映射的性质表明，对任何 $X\in\mathrm{M}_n(\mathbb R)$，映射 $t\in\mathbb R\mapsto\exp(tX)\in\mathrm{GL}_n(\mathbb R)$ 是一个单参数子群. 事实上，$\mathrm{GL}_n(\mathbb R)$ 的所有单参数子群都具有这种形式. 更确切地说：
\begin{proposition*}
设 $\lambda:\mathbb R\longrightarrow\mathrm{GL}_n(\mathbb R)$ 是一个单参数子群，并设 $\lambda'(0)\in\mathrm{M}_n(\mathbb R)$ 为其在 $0$ 处的导数. 那么对一切 $t\in\mathbb R$，有 $\lambda(t)=\exp(t\lambda'(0))$.
\end{proposition*}
\begin{proof}
由于 $\lambda$ 是群同态，对一切 $t\in\mathbb R$ 有 $\lambda'(t)=\lambda(t)\lambda'(0)$. 同样地，对一切 $t$ 有
\[
\frac{d}{dt}\bigl(t\mapsto\exp(t\lambda'(0))\bigr)=\exp(t\lambda'(0))\lambda'(0).
\]
于是 $\lambda$ 与 $t\mapsto\exp(t\lambda'(0))$ 满足同一个线性微分方程以及同一个初始条件 $\lambda(0)=\exp(0\lambda'(0))=e$. 由这类方程解的唯一性，这两个函数在整个 $\mathbb R$ 上相等.
\end{proof}
我们记 $\mathscr X(G)$ 为 Lie 群 $G$ 的单参数子群全体构成的集合. 上面的命题表明，映射
\[
\mathscr X(\mathrm{GL}_n(\mathbb R))\longrightarrow\mathfrak{gl}_n(\mathbb R),\ \lambda\mapsto\lambda'(0)
\]
是一个\emph{双射}，其逆为 $X\mapsto(t\mapsto\exp(tX))$. 于是人们自然会问：如何用单参数子群的语言来表述 $\mathfrak{gl}_n(\mathbb R)$ 的向量结构与其 Lie 括号.
\end{numpara}
\begin{proposition}\label{prop:lie-1-6-3}
$\mathscr X(\mathrm{GL}_n(\mathbb R))$ 上的 $\mathbb R$-Lie 代数结构由下列公式给出：
\begin{itemize}[label=--]
\item $(r.\lambda)(t)=\lambda(rt)$，其中 $r\in\mathbb R$；
\item $(\lambda+\mu)(t)=\lim_{k\to\infty}(\lambda(t/k)\mu(t/k))^k$；
\item $[\lambda,\mu](t)=\lim_{k\to\infty}(\lambda(1/k)\mu(t/k)\lambda(-1/k)\mu(-t/k))^{k^2}$.
\end{itemize}
\end{proposition}
\begin{proof}
关于乘以 $r$ 的公式是显然的. 对于加法与括号，只需证明下列两个公式：
\begin{itemize}[label=--]
\item $\lim_{k\to+\infty}\bigl(\exp(\tfrac{X}{k})\exp(\tfrac{Y}{k})\bigr)^{k}=\exp(X+Y)$；
\item $\lim_{k\to+\infty}\bigl(\exp(\tfrac{X}{k})\exp(\tfrac{Y}{k})(\exp(\tfrac{-X}{k})\exp(\tfrac{-Y}{k}))\bigr)^{k^2}=\exp([X,Y])$.
\end{itemize}
这两个公式只需作有限阶展开即可得到，\emph{参见}本人 M1 课程讲义 2.1.2.
\end{proof}
\begin{theorem}\label{thm:lie-1-6-4}
设 $G$ 是 $\mathrm{GL}_n(\mathbb R)$ 的一个闭子群. 那么 $G$ 是一个 Lie 子群，并且
\[
\mathrm{Lie}(G)=\{X\in\mathrm{M}_n(\mathbb R),\ \forall t\in\mathbb R,\ \exp(tX)\in G\}.
\]
\end{theorem}
\begin{proof}
令 $L:=\{X\in\mathrm{M}_n(\mathbb R),\ \forall t\in\mathbb R,\ \exp(tX)\in G\}$. 关键的一点是，根据下面的命题，这个集合是一个 $\mathbb R$-子向量空间\footnote{译注：原文如此（``d'après la proposition suivante''，即``根据下面的命题''），但本节中该处之后并无紧随其后的命题.}.
在 $\mathrm{M}_n(\mathbb R)$ 中取 $L$ 的一个补子空间 $E$，并考虑映射
\[
F:L\times E\longrightarrow\mathrm{GL}_n(\mathbb R),\ (X,Y)\mapsto\exp(X)\exp(Y).
\]
其微分 $d_{(0,0)}F$ 可逆（它就是求和映射 $L\oplus E\longrightarrow\mathrm{M}_n(\mathbb R)$）. 因此存在
\begin{itemize}[label=--]
\item $L$ 中 $0$ 的一个邻域 $\mathcal U$，
\item $E$ 中 $0$ 的一个邻域 $\mathcal V$，
\item 以及 $\mathrm{GL}_n(\mathbb R)$ 中 $I_n$ 的一个邻域 $\mathcal W$，
\end{itemize}
使得 $F$ 实现微分同胚 $\mathcal U\times\mathcal V\stackrel{\sim}{\longrightarrow}\mathcal W$. 注意 $F(\mathcal U\times\{0\})\subset\mathcal W\cap G$.
我们将证明：存在 $E$ 中 $0$ 的一个邻域 $\mathcal V'\subset\mathcal V$，使得令 $\mathcal W'=F(\mathcal U\times\mathcal V')$ 时，有
\[
(*)\ :\ F(\mathcal U\times\{0\})=\mathcal W'\cap G.
\]
这将表明 $\mathcal W'\cap G$ 是 $\mathcal W'$ 的一个子流形，并且把这个子流形在 $e$ 处的切空间等同于 $\mathcal U$ 在 $0$ 处的切空间，即 $L$. 再经过平移，最终可推得 $G$ 是 $\mathrm{GL}_n(\mathbb R)$ 的子流形，且 $T_eG=L$，这正是所要证的.
我们先暂且承认
\[
(**)\ :\ \text{存在 $E$ 中 $0$ 的一个邻域 $\mathcal V'\subset\mathcal V$，使得 $\exp(\mathcal V')\cap G=\{I_n\}$.}
\]
于是令 $\mathcal W'=F(\mathcal U\times\mathcal V')$. 对任何 $g\in\mathcal W'\cap G$，存在 $X\in\mathcal U$，$Y\in\mathcal V'$，使得 $g=F(X,Y)$. 但此时 $\exp(Y)=\exp(-X)g\in\exp(\mathcal V')\cap G=\{I_n\}$，故 $g=\exp(X)$. 由此可知等式 $(*)$ 对 $\mathcal V'$ 成立.
剩下要证明 $(**)$. 假设不然. 那么可以找到 $E$ 中一列\emph{非零}元素 $(Y_n)_{n\in\mathbb N}$ 趋于 $0$，使得对一切 $n$ 有 $\exp(Y_n)\in G$. 在抽取一个子列之后，还可以假设序列 $(Y_n/\|Y_n\|)_{n\in\mathbb N}$ 在 $E$ 中有极限 $Y$（球面的紧性）. 注意 $Y\neq0$. 我们将证明 $Y\in L$，这将与 $L\cap E=\{0\}$ 矛盾. 为此我们需要证明：对一切 $t\in\mathbb R$ 有 $\exp(tY)\in G$. 记 $[x]$ 为实数 $x$ 的整数部分，并写出
\[
\exp(t/\|Y_n\|\cdot Y_n)=\exp(Y_n)^{[t/\|Y_n\|]}\cdot\exp((t/\|Y_n\|- [t/\|Y_n\|])Y_n).
\]
由于项 $\exp((t/\|Y_n\|- [t/\|Y_n\|])Y_n)$ 当 $n$ 趋于 $+\infty$ 时趋于 $1$，我们看到
\[
\exp(tY)=\lim_{n\to+\infty}\exp(Y_n)^{[t/\|Y_n\|]}\in G.
\]
\end{proof}
\begin{numpara}[经典 Lie 代数]\label{np:lie-1-6-5}
我们现在知道经典群的确都是 Lie 群，并且有了计算其 Lie 代数的方法. 所有经典群都定义在某个适当的 $\mathrm{GL}_n(K)$ 或 $\mathrm{GL}(V)$ 之中，其 Lie 代数的括号由 $\mathfrak{gl}_n$ 或 $\mathfrak{gl}(V)$ 的括号诱导. 因此我们只需描述作为基础的 $\mathbb R$-向量空间所满足的定义条件.
例如，等式 $\exp(\operatorname{tr}(u))=\det(\exp(u))$ 表明\footnote{译注：原书下式中 $\mathrm{SL}(V)$ 之后缺一个右括号，应为 $\mathrm{Lie}(\mathrm{SL}(V))=\{\cdots\}$.}
\[
\mathfrak{sl}(V):=\mathrm{Lie}(\mathrm{SL}(V)=\{u\in\mathfrak{gl}(V),\ \operatorname{tr}(u)=0\}.
\]
注意它总是一个 $K$-向量空间（而不仅仅是一个 $\mathbb R$-向量空间）. 用矩阵的语言，这给出
\begin{enumerate}[label=\roman*)]
\item $\mathfrak{sl}_n(\mathbb R)=\{X\in\mathfrak{gl}_n(\mathbb R),\ \operatorname{tr}(X)=0\}$（迹为零的矩阵）. 其维数为 $n^2-1$.
\item $\mathfrak{sl}_n(\mathbb C)=\{X\in\mathfrak{gl}_n(\mathbb C),\ \operatorname{tr}(X)=0\}$. 其实维数为 $2n^2-2$.
\end{enumerate}
另一方面，设 $G=\{g\in\mathrm{GL}(V),\ \Psi(gv,gw),\ \forall v,w\in V\}$\footnote{译注：原书此处排版有脱漏，按上下文应为 $\Psi(gv,gw)=\Psi(v,w)$.}由一个非退化双线性型 $\Psi$（无论它是对称的、交错的还是埃尔米特的）定义. 换言之，若 $g^*$ 表示 $g$ 关于 $\Psi$ 的伴随，则有 $G=\{g\in\mathrm{GL}(V),\ gg^*=\operatorname{id}_V\}$. 由于 $\exp(u)^*=\exp(u^*)$（通过取极限），我们得到
\[
\mathfrak g=\{u\in\mathfrak{gl}(V),\ u+u^*=0\}.
\]
顺便指出，它总是一个 $K$-向量空间，唯独在酉的情形例外：此时 $K=\mathbb C$，但 $\mathfrak g$ 只是一个 $\mathbb R$-向量空间. 用矩阵的语言，这给出：
\begin{enumerate}[label=\roman*),start=3]
\item $\mathfrak{so}_n(\mathbb R)=\mathrm{Lie}(\mathrm{O}(n))=\mathrm{Lie}(\mathrm{SO}(n))=\{X\in\mathfrak{gl}_n(\mathbb R),\ X+{}^tX=0\}$，即反对称矩阵全体. 其维数为 $n(n-1)/2$.
\item $\mathfrak{so}_n(\mathbb C)=\mathrm{Lie}(\mathrm{O}(n,\mathbb C))=\{X\in\mathfrak{gl}_n(\mathbb C),\ X+{}^tX=0\}$. 其实维数为 $n(n-1)$.
\item $\mathfrak u_n=\mathrm{Lie}(\mathrm{U}(n))=\{X\in\mathfrak{gl}_n(\mathbb C),\ X+{}^t\overline{X}=0\}$. 例如，$\mathfrak u_1=i\mathbb R\subset\mathfrak{gl}_1(\mathbb C)=\mathbb C$. 其维数为 $n^2$.
\item $\mathfrak{su}_n=\mathrm{Lie}(\mathrm{SU}(n))=\mathfrak u_n\cap\mathfrak{sl}_n(\mathbb C)$. 其维数为 $n^2-1$.
\item $\mathfrak{sp}_{2n}(K)=\mathrm{Lie}(\mathrm{Sp}_{2n}(K))=\{X\in\mathfrak{gl}_{2n}(K),\ {}^tXJ_{2n}+J_{2n}X=0\}$. 为计算维数，我们注意到 ${}^tXJ_{2n}+J_{2n}X=0\Leftrightarrow$（$J_{2n}X$ 反对称）. 我们求得：当 $K=\mathbb R$ 时实维数为 $n(n+1)/2$，当 $K=\mathbb C$ 时为 $n(n+1)$.
\end{enumerate}
\end{numpara}
\begin{numpara}[小结]\label{np:lie-1-6-6}
对于 $\mathrm{GL}_n(\mathbb R)$ 的闭子群 $G$，我们现在知道：
\begin{enumerate}[label=\roman*)]
\item $G$ 是一个 Lie 群.
\item 映射 $\mathscr X(G)\longrightarrow\mathrm{Lie}(G)$，$\lambda\mapsto\lambda'(0)$ 是双射.
\item $\mathscr X(G)$ 上的 $\mathbb R$-Lie 代数结构由命题~\ref{prop:lie-1-6-3} 的公式给出.
\item 映射 $\exp_G:\mathrm{Lie}(G)\simeq\mathscr X(G)\longrightarrow G$，$\lambda\mapsto\lambda(1)$ 实现了从 $\mathrm{Lie}(G)$ 中 $0$ 的一个邻域到 $G$ 中 $e$ 的一个邻域的微分同胚.
\item 若 $\varphi:G\longrightarrow G'$ 是到另一个 $\mathrm{GL}_{n'}(\mathbb R)$ 的闭子群的 Lie 群同态，则 $d\varphi$ 可等同于复合映射 $\mathscr X(G)\longrightarrow\mathscr X(G')$，$\lambda\mapsto\varphi\circ\lambda$. 特别地，有 $\exp_{G'}\circ d\varphi=\varphi\circ\exp_G$.
\end{enumerate}
\noindent\textbf{例.}
\begin{enumerate}[label=\roman*)]
\item 取 $\varphi=\operatorname{Int}_x$，得到 $\operatorname{Int}_x\circ\exp_G=\exp_G\circ\operatorname{Ad}_x$.
\item 取 $\varphi=\operatorname{Ad}:G\longrightarrow\mathrm{GL}(\mathrm{Lie}(G))$，得到：对一切 $X\in\mathrm{Lie}(G)$，$\operatorname{Ad}_{\exp_G(X)}=\exp_{\mathrm{GL}(T_eG)}(\operatorname{ad}_X)$. 即便在 $G=\mathrm{GL}_n(\mathbb R)$ 的情形，矩阵等式（对 $Y\in\mathrm{M}_n(\mathbb R)$）$\exp(X)Y\exp(-X)=\exp(\operatorname{ad}_X)(Y)$ 也难以直接证明.
\item 若 $\varphi$ 是一个嵌入 $G\hookrightarrow\mathrm{GL}_n(\mathbb R)$，则 $d\varphi$ 是一个嵌入 $\mathrm{Lie}(G)\hookrightarrow\mathfrak{gl}_n(\mathbb R)$，且 $G$ 的指数映射是 $\mathrm{GL}_n(\mathbb R)$ 中矩阵指数映射的限制.
\end{enumerate}
\end{numpara}
请注意，这些性质的陈述并未用到 $G$ 到某个 $\mathrm{GL}_n(\mathbb R)$ 中的嵌入（在证明中只用到了这种嵌入的存在性）. 因此我们有希望把它们推广到任意 Lie 群，这正是下一节要做的事.

\emph{术语}：容许一个到某个 $\mathrm{GL}_n(\mathbb R)$ 中的嵌入的 Lie 群称为\newterm{lie-linear-lie-group}{线性 Lie 群}{linear Lie group}.
下面给出这些性质的一个推论的例子：
\begin{proposition}\label{prop:lie-1-6-7}
设 $G$ 是一个连通的线性 Lie 群. 伴随表示 $\operatorname{Ad}_G:G\longrightarrow\mathrm{GL}(\mathrm{Lie}(G))$ 的核是 $G$ 的中心 $Z(G)$.
\end{proposition}
\begin{proof}
$Z(G)\subset\operatorname{Ker}(\operatorname{Ad}_G)$ 这一事实是容易的：若 $z\in Z(G)$，则 $\operatorname{Int}_z=\operatorname{id}_G$，从而 $\operatorname{Ad}_z=\operatorname{id}_{T_eG}$. 反之，设 $x\in\operatorname{Ker}(\operatorname{Ad}_G)$. 根据小结~\ref{np:lie-1-6-6}~中的性质 v)，有 $\exp_G\circ\operatorname{Ad}_x=\operatorname{Int}_x\circ\exp_G$. 再根据性质 iv) 可知，$\operatorname{Int}_x$ 在 $e$ 的某个邻域 $\mathcal U$ 上恒等作用. 由于被 $\operatorname{Int}_x$ 固定的 $G$ 中元素全体构成一个群，它包含由邻域 $\mathcal U$ 生成的群，即（由于 $G$ 连通）$G$ 本身. 故 $\operatorname{Int}_x=\operatorname{id}_G$，从而 $x\in Z(G)$.
\end{proof}
\begin{exercise*}
我们要证明 $\mathrm{SU}(2)$ 是 $\mathrm{SO}(3,\mathbb R)$ 的万有覆盖.
\begin{enumerate}[label=\roman*)]
\item 证明 $\mathrm{SU}(2)$ 在球面 $S^3$ 上单可递地作用，并由此推出它是单连通的.
\item 证明 Lie 代数 $\mathfrak{su}(2)$ 以下列矩阵为一组基：
\[
\begin{pmatrix}0&1\\-1&0\end{pmatrix},\
\begin{pmatrix}0&i\\i&0\end{pmatrix},\
\text{以及}\ \begin{pmatrix}i&0\\0&-i\end{pmatrix}.
\]
\item 考虑表示 $\operatorname{Ad}_{\mathrm{SU}(2)}:\mathrm{SU}(2)\longrightarrow\mathrm{GL}(\mathfrak{su}(2))$. 借助上述基把 $\mathrm{GL}(\mathfrak{su}(2))$ 等同于 $\mathrm{GL}_3(\mathbb R)$. 然后证明像 $\operatorname{Ad}(\mathrm{SU}(2))$ 含于 $\mathrm{SO}(3,\mathbb R)\subset\mathrm{GL}_3(\mathbb R)$.
\item 证明 $\operatorname{Ker}\operatorname{Ad}_{\mathrm{SU}(2)}=\{\pm I_2\}$. 由此推出 $\operatorname{ad}_{\mathfrak{su}(2)}$ 诱导同构 $\mathfrak{su}(2)\stackrel{\sim}{\longrightarrow}\mathfrak{so}(3,\mathbb R)$，进而 $\operatorname{Ad}_{\mathrm{SU}(2)}$ 诱导同构 $\mathrm{SU}(2)/\{\pm I_2\}\stackrel{\sim}{\longrightarrow}\mathrm{SO}(3,\mathbb R)$.
\end{enumerate}
\end{exercise*}
关于从四元数 $\mathbb H$ 出发对同态 $\mathrm{SU}(2)\longrightarrow\mathrm{SO}(3,\mathbb R)$ 的另一种构造，\emph{参见} M1 讲义 1.3.6.
